\section{Introduction}
\label{secIntroductionChaine}

Le premier TP a permis de mettre en service la communication avec le
contrôleur \textbf{PFC200} (poste guichet) : chaque bit produit ou
consommé par cet équipement correspond directement à un capteur ou un
actionneur, et un simple \textbf{alias mémoire} (chevauchement d'une
variable structurée sur la zone \texttt{\%MW} du coupleur) suffit à lui
donner un nom exploitable dans le programme.

Ce document prépare la suite des TP, où deux autres équipements du réseau
\texttt{NOC\_ETHIP\_ROBOT} doivent être mis en œuvre : le contrôleur
\textbf{M221} (pupitre opérateur et écran de supervision) et la baie
\textbf{CS9} du robot (serveur \texttt{uniVALplc}). Dans les deux cas, un
alias seul ne suffit plus : les octets échangés ne portent pas
directement une information unitaire exploitable, il faut un
\textbf{bloc fonctionnel de conversion} pour leur donner du sens.

\begin{UPSTIinfor}{Pourquoi un alias ne suffit plus}
Pour le PFC200, chaque bit \emph{est} l'information (un capteur, une
commande). Pour le M221 et pour la baie CS9, les octets échangés sont le
résultat d'un \textbf{encodage} propre à la bibliothèque logicielle de
l'équipement (empaquetage de plusieurs informations dans un même octet,
structures imbriquées, types numériques codés sur plusieurs octets...).
Un alias donne seulement accès à cette zone \textbf{brute} : il faut un
bloc fonctionnel qui la décode (en lecture) ou qui l'encode (en écriture)
pour manipuler des informations ayant un nom et un type clairs dans le
programme.
\end{UPSTIinfor}

Ce document traite cette chaîne de conversion dans l'ordre suivant :

\begin{enumerate}
    \item le contrôleur \textbf{M221} (blocs \texttt{FB\_ReadHmi} et
    \texttt{FB\_WriteHmi}), le cas le plus simple des deux ;
    \item la baie \textbf{CS9} (blocs \texttt{VAL\_ReadAxesGroup} et
    \texttt{VAL\_WriteAxesGroup}), plus riche ;
    \item un schéma de synthèse comparant les trois chaînes de conversion
    rencontrées (PFC200, M221, CS9) ;
    \item un TD d'application, sur la mise en œuvre du mode GEMMA
    \textbf{A6} (mise en conditions initiales) du robot.
\end{enumerate}
